\chapter{Implementation Gap Review and Agent Rule Pack}
\label{chap:simula-gap-review}

\section{Review scope and constraints}
This chapter records a repository review performed against the
\emph{Simula Training Data Framework} specification PDF and the
current code under \texttt{/Users/user/Documents/sap-oss/src} as of
2026-04-21. The review covered four targets:
\begin{enumerate}[label=\arabic*.]
  \item the normative Simula implementation path in
    \texttt{src/training/};
  \item the implementation-agent and maintenance rule location under
    \texttt{src/training/};
  \item the operational frontend and API wrapper in
    \texttt{src/generativeUI/training-webcomponents-ngx/};
  \item runtime monitoring and alerting support for Simula runs.
\end{enumerate}

Two delivery constraints shaped the outcome:
\begin{itemize}
  \item This workstream was instructed to change only files under
    \texttt{docs/}. No source changes were made under
    \texttt{src/}.
  \item The intended agent location is \texttt{src/training/}. Because
    source-tree edits were out of scope for this task, the
    implementation-agent and runtime-monitor definitions are captured
    here as canonical proposed artefacts to be materialised later
    under \texttt{src/training/.clinerules} and a companion runtime
    monitor rule file in the same subtree.
\end{itemize}

\section{Executive assessment}
\begin{table}[htbp]
\centering
\footnotesize
\caption{Current implementation status against the Simula specification.}
\label{tab:simula-gap-summary}
\begin{tabularx}{\textwidth}{@{}l l X@{}}
\toprule
\textbf{Area} & \textbf{Status} & \textbf{Assessment} \\
\midrule
Default CLI and pipeline entrypoints & \textbf{Red} & The repository default still executes the legacy CSV/template/Spider pipeline, not the HANA-direct Simula flow described in chapters~\ref{chap:extraction}--\ref{chap:generator}. \\
Simula core modules & \textbf{Amber} & HANA extraction, taxonomy building, generation, and evaluator modules exist, but they are only partially integrated and several evaluation components fail at runtime because their data contracts drifted. \\
Evaluation framework & \textbf{Red} & Complexity calibration, coverage, diversity, taxonomy quality, and critic calibration are present as isolated modules, but they are not orchestrated by the runtime path that produces datasets. \\
Frontend training console & \textbf{Amber/Red} & The shell, governance surfaces, and HANA explorer exist, but the pipeline UX still models the older workflow and does not expose Simula-specific artefacts or evaluator outputs. \\
Runtime monitoring and alerting & \textbf{Amber/Red} & Generic health and notification primitives exist, but there is no Simula-specific runtime monitor that watches spec-critical outputs, thresholds, and evaluator failures. \\
\texttt{src/training} agent rule materialisation & \textbf{Amber} & The canonical agent location is \texttt{src/training}, but the rule files are documented here rather than created in source because this task was constrained to \texttt{docs/}. \\
\bottomrule
\end{tabularx}
\end{table}

\section{Validation performed}
The review did not rely only on static reading. The following checks
were run against the current repository state:
\begin{itemize}
  \item \texttt{pytest -q src/training/pipeline/tests/test\_pipeline.py}
    passed (\textbf{20/20}); this confirms that the \emph{legacy}
    pipeline remains internally healthy.
  \item \texttt{pytest -q packages/api-server/tests/test\_main.py -k
    'hana\_stats\_falls\_back\_to\_preview\_when\_unconfigured or
    hana\_query\_returns\_preview\_rows\_when\_unconfigured'} passed
    (\textbf{2/2}); this confirms that preview fallback behaviour is
    functioning as designed in the current API wrapper.
  \item Runtime reproductions demonstrated that
    \texttt{SimulaCriticValidator} and
    \texttt{SimulaCoverageEvaluator} are not execution-safe against
    the current \texttt{TrainingExample} and \texttt{Taxonomy}
    classes: the observed failures were
    \texttt{TypeError: TrainingExample.\_\_init\_\_() got an
    unexpected keyword argument 'schema\_context'} and
    \texttt{AttributeError: 'Taxonomy' object has no attribute 'id'}.
\end{itemize}

\section{Findings in \texttt{src/training}}

\subsection*{SIM-GAP-001 (P0): the default execution path is still the legacy pipeline}
\textbf{Observed state.}
\texttt{src/training/pipeline/main.py} still exposes the
legacy command set
(\texttt{extract-schema}, \texttt{parse-templates},
\texttt{expand}, \texttt{format-spider}, \texttt{upload}),
and the default Makefile target still executes the same
CSV/template/Spider flow.

\textbf{Evidence.}
\begin{itemize}
  \item \texttt{src/training/pipeline/main.py:3--13, 30--39,
    136--141}
  \item \texttt{src/training/pipeline/Makefile:4--11, 29--30,
    35--82}
  \item By contrast, the Simula orchestration path lives in
    \texttt{src/training/scripts/run\_hana\_pipeline.py:3--10,
    51--196} and is not the repository default.
\end{itemize}

\textbf{Impact.}
An operator following the current default entrypoint runs the old
pipeline, not the spec-defined HANA extraction, taxonomy generation,
and reasoning-driven synthesis flow. This is the single largest
end-to-end compliance gap.

\textbf{Required remediation.}
Promote the Simula path to the authoritative CLI/runtime entrypoint,
and demote or quarantine the legacy template pipeline behind an
explicit compatibility switch.

\subsection*{SIM-GAP-002 (P0): API and job orchestration are wired to the legacy generators}
\textbf{Observed state.}
The training console backend does not call
\texttt{run\_hana\_pipeline.py}. Instead:
\begin{itemize}
  \item the pipeline websocket worker runs \texttt{make all} in
    \texttt{src/training/pipeline/};
  \item the training-generation job runner executes
    \texttt{schema\_pipeline/data\_generator.py} followed by
    \texttt{prepare\_training\_data.py}.
\end{itemize}

\textbf{Evidence.}
\begin{itemize}
  \item \texttt{packages/api-server/src/main.py:2904--3019}
  \item \texttt{packages/api-server/src/main.py:2222--2357}
\end{itemize}

\textbf{Impact.}
Even if the Simula modules were complete, the deployed frontend/API
stack would still launch the older data path and publish the older
artefact contract. The Simula implementation therefore remains
present in code but absent from the operational surface.

\textbf{Required remediation.}
Replace the API worker commands with the Simula runner and its report
artefacts; keep the current legacy commands only for rollback or
migration support.

\subsection*{SIM-GAP-003 (P0): evaluation modules have drifted away from the current data contracts}
\textbf{Observed state.}
The current \texttt{TrainingExample} dataclass contains only
\texttt{id}, \texttt{question}, \texttt{sql}, \texttt{domain},
\texttt{table}, \texttt{difficulty}, \texttt{taxonomy\_mix},
\texttt{meta\_prompt}, \texttt{is\_complexified}, and
\texttt{critic\_verdict}. However:
\begin{itemize}
  \item \texttt{SimulaCriticValidator} constructs
    \texttt{TrainingExample} objects using nonexistent fields such as
    \texttt{schema\_context}, \texttt{taxonomy\_path},
    \texttt{complexity\_score}, \texttt{critic\_passed}, and
    \texttt{generation\_metadata};
  \item \texttt{SimulaCoverageEvaluator} expects each
    \texttt{Taxonomy} to expose \texttt{id} and \texttt{factor}, and
    each \texttt{TaxonomyNode} to expose \texttt{id}; the current
    taxonomy classes expose \texttt{factor\_name} and
    \texttt{name}/\texttt{description}/\texttt{children} instead.
\end{itemize}

\textbf{Evidence.}
\begin{itemize}
  \item \texttt{src/training/pipeline/simula\_data\_generator.py:32--44}
  \item \texttt{src/training/pipeline/simula\_critic\_validator.py:212--238, 268--272}
  \item \texttt{src/training/pipeline/simula\_coverage\_evaluator.py:157--170, 203--215}
  \item Runtime reproductions raised the concrete exceptions recorded
    in the validation section above.
\end{itemize}

\textbf{Impact.}
The spec's evaluation chapter is not merely unintegrated; two of its
core evaluators are currently unusable against the live Simula data
model. This blocks critic calibration, coverage measurement, and any
automatic gap-filling workflow built on top of them.

\textbf{Required remediation.}
Unify the \texttt{TrainingExample}, \texttt{Taxonomy}, and evaluator
contracts before further orchestration work. Until those contracts
are repaired, downstream evaluation integration should be treated as
blocked.

\subsection*{SIM-GAP-004 (P1): the Simula orchestration script stops before the full evaluation framework}
\textbf{Observed state.}
\texttt{run\_hana\_pipeline.py} performs three phases:
\begin{enumerate}[label=\arabic*.]
  \item HANA schema extraction;
  \item taxonomy construction;
  \item example generation and file write-out.
\end{enumerate}
It writes \texttt{extracted\_schema.json},
\texttt{taxonomies.json}, a timestamped training JSONL, and
\texttt{pipeline\_stats.json}. It does \emph{not} call the evaluator
modules for complexity calibration, coverage, diversity, taxonomy
quality, critic calibration, or adaptive gap filling.

\textbf{Evidence.}
\begin{itemize}
  \item \texttt{src/training/scripts/run\_hana\_pipeline.py:51--196}
  \item Search results show that
    \texttt{SimulaComplexityCalibrator},
    \texttt{SimulaCoverageEvaluator},
    \texttt{SimulaCriticValidator},
    \texttt{SimulaDiversityAnalyzer},
    \texttt{SimulaDiversityOptimizer}, and
    \texttt{SimulaTaxonomyEvaluator} are defined and re-exported, but
    not invoked from any orchestration path.
\end{itemize}

\textbf{Impact.}
The implementation currently reaches chapter~5
(\emph{Generation Engine}) but not chapter~7
(\emph{Evaluation Framework}) in an operational sense. The spec's
quality gates therefore exist largely as library code and prose, not
as emitted run artefacts.

\textbf{Required remediation.}
Extend the Simula runner so that every completed generation run emits
all evaluation reports required by the spec and records their
threshold outcomes into governance and monitoring surfaces.

\subsection*{SIM-GAP-005 (P1): adaptive controls are specified in configuration but not applied in the generator}
\textbf{Observed state.}
\texttt{GenerationConfig} defines
\texttt{get\_adaptive\_complexity\_ratio()} and
\texttt{get\_adaptive\_critic\_threshold()}, but the main generation
loop continues to use the raw \texttt{complexity\_ratio} and a plain
double-critic pass/fail call without consulting those helpers.

\textbf{Evidence.}
\begin{itemize}
  \item \texttt{src/training/pipeline/simula\_config.py:269--326}
  \item \texttt{src/training/pipeline/simula\_data\_generator.py:494--500, 599--614}
  \item Repository-wide search found no call sites for
    \texttt{get\_adaptive\_complexity\_ratio()} or
    \texttt{get\_adaptive\_critic\_threshold()} outside their
    definitions.
\end{itemize}

\textbf{Impact.}
Adaptive features introduced in the v1.2 spec are represented as
configuration helpers only. The runtime path does not yet change its
behaviour in response to teacher strength or measured complexity.

\textbf{Required remediation.}
Thread the adaptive controls through generation, critic acceptance,
and report emission so the pipeline can prove that the v1.2 adaptive
logic is active.

\subsection*{SIM-GAP-006 (P1): output contracts do not line up across the Simula runner, API, and UI}
\textbf{Observed state.}
The Simula runner writes timestamped files such as
\texttt{simula\_training\_\{timestamp\}.jsonl} under
\texttt{data/hana\_generated/}. The API and frontend continue to
expect the older path \texttt{training/pipeline/output/train.jsonl}.

\textbf{Evidence.}
\begin{itemize}
  \item \texttt{src/training/pipeline/simula\_data\_generator.py:672--716}
  \item \texttt{packages/api-server/src/main.py:2960--2963,
    3007--3012, 3354--3360}
  \item \texttt{apps/angular-shell/src/app/pages/pipeline/pipeline.component.ts:549--554}
\end{itemize}

\textbf{Impact.}
Even after the Simula runner is wired in, the frontend and API will
not discover the artefacts they are looking for unless the output
contract is normalised.

\textbf{Required remediation.}
Choose one authoritative artefact layout for Simula runs and apply it
consistently across runner, backend, governance records, preview
endpoints, and UI components.

\subsection*{SIM-GAP-007 (P2): automated tests still prove only the legacy pipeline path}
\textbf{Observed state.}
The repository contains passing unit tests for the legacy pipeline
modules, but there are no corresponding tests for
\texttt{run\_hana\_pipeline.py} or the end-to-end Simula evaluator
stack.

\textbf{Evidence.}
\begin{itemize}
  \item \texttt{src/training/pipeline/tests/test\_pipeline.py}
  \item The passing test run recorded in the validation section above
    exercised only those legacy tests.
\end{itemize}

\textbf{Impact.}
The project currently has high confidence in the deprecated path and
low confidence in the normative one.

\textbf{Required remediation.}
Add unit and integration coverage for extraction, taxonomy, critic
filtering, evaluator compatibility, output contracts, and the API
worker that launches the Simula flow.

\section{Findings in \texttt{src/generativeUI/training-webcomponents-ngx}}

\subsection*{UI-GAP-001 (P0): the pipeline UI still models the old six-stage workflow}
\textbf{Observed state.}
The pipeline page shows the following stages:
\texttt{Preconvert}, \texttt{Extract Schema}, \texttt{Parse
Templates}, \texttt{Expand}, \texttt{Build SQL}, and
\texttt{Format}. Stage transitions are inferred from keywords such as
\texttt{converting}, \texttt{parsing}, and \texttt{formatting}.

\textbf{Evidence.}
\begin{itemize}
  \item \texttt{apps/angular-shell/src/app/pages/pipeline/pipeline.component.ts:333--339}
  \item \texttt{apps/angular-shell/src/app/pages/pipeline/pipeline.component.ts:488--497}
\end{itemize}

\textbf{Impact.}
The operational UI tells users that the product pipeline is still the
preconvert/template/Spider flow. It does not surface the Simula phase
model from this spec: HANA extraction, taxonomy building, data
generation, and evaluation.

\textbf{Required remediation.}
Replace the stage model with the Simula phase graph and emit
stage-specific artefacts and metrics from the backend.

\subsection*{UI-GAP-002 (P1): runtime telemetry in the pipeline page is simulated rather than sourced from the backend}
\textbf{Observed state.}
The pipeline page's worker grid and throughput readout are generated
from random numbers whenever the pipeline state is \texttt{running}.

\textbf{Evidence.}
\begin{itemize}
  \item \texttt{apps/angular-shell/src/app/pages/pipeline/pipeline.component.ts:369--379}
\end{itemize}

\textbf{Impact.}
The page provides motion, but not observability. It cannot support a
runtime-monitoring role because the visual telemetry is not derived
from actual execution state.

\textbf{Required remediation.}
Bind the concurrency, throughput, and stage timing displays to real
worker metrics, log events, and evaluator outcomes.

\subsection*{UI-GAP-003 (P1): the frontend/API operational surface is still wired to preview or compatibility paths}
\textbf{Observed state.}
The HANA explorer and related backend endpoints deliberately fall back
to preview data when credentials or connectivity are missing. The API
wrapper exposes this cleanly, but the Simula pipeline has no explicit
distinction between a live production run and a preview-backed
demonstration surface.

\textbf{Evidence.}
\begin{itemize}
  \item \texttt{packages/api-server/src/main.py:3030--3038,
    3149--3197, 3349--3395}
  \item \texttt{apps/angular-shell/src/app/pages/hana-explorer/hana-explorer.component.ts:17--25, 73--82, 112--118}
  \item API preview tests passed, confirming that preview mode is an
    intended and working behaviour in the current stack.
\end{itemize}

\textbf{Impact.}
Preview fallback is appropriate for resilience, but the current
product surface does not yet enforce the sharper contract the Simula
spec needs: production corpus generation must be visibly live, never
silently satisfied by preview artefacts.

\textbf{Required remediation.}
Add explicit live-vs-preview run state, block governed production
launches when required dependencies are only in preview mode, and
surface that policy directly in the UI.

\subsection*{UI-GAP-004 (P1): there is no Simula-specific UI for taxonomy and evaluation artefacts}
\textbf{Observed state.}
The frontend contains generic dashboard, analytics, registry,
governance, HANA explorer, and data-product screens, but there is no
route or component dedicated to:
\begin{itemize}
  \item extracted schema inspection as a Simula run artefact;
  \item taxonomy visualisation and taxonomy quality review;
  \item coverage, diversity, critic calibration, or Elo complexity
    reports;
  \item student-teacher gap tracking or adaptive-complexity
    decisions.
\end{itemize}

\textbf{Evidence.}
\begin{itemize}
  \item \texttt{apps/angular-shell/src/app/app.routes.ts:1--92}
  \item Repository-wide searches for
    \texttt{SimulaCoverageEvaluator},
    \texttt{SimulaCriticValidator},
    \texttt{SimulaComplexityCalibrator},
    \texttt{SimulaDiversityAnalyzer}, and
    \texttt{SimulaTaxonomyEvaluator} found no corresponding Angular
    surfaces.
\end{itemize}

\textbf{Impact.}
The spec's evaluator outputs cannot yet be reviewed in the operational
console, which means human sign-off, explainability, and runtime
triage still depend on raw files or logs rather than first-class
product screens.

\textbf{Required remediation.}
Add explicit Simula result views or extend existing governance and
registry pages so that each spec-defined artefact becomes inspectable.

\subsection*{UI-GAP-005 (P2): notification and alerting are generic, not Simula-specific}
\textbf{Observed state.}
Notification CRUD exists in the API, but the shell notification tray
is currently composed from generic health, pipeline state, and GPU
pressure held in the store. No Simula-specific alert rules are
registered for evaluator failures, missing artefacts, contract drift,
or unexpected preview mode.

\textbf{Evidence.}
\begin{itemize}
  \item \texttt{packages/api-server/src/notification\_api.py}
  \item \texttt{packages/api-server/src/main.py:590}
  \item \texttt{apps/angular-shell/src/app/components/shell/shell.component.ts:645--691}
\end{itemize}

\textbf{Impact.}
The plumbing for notifications exists, but the actual runtime monitor
requested for Simula has not yet been encoded as alerting logic.

\textbf{Required remediation.}
Back the notification surface with rules derived from Simula run
artefacts and thresholds, not only generic platform health.

\section{Proposed implementation-agent rule file}
\label{sec:simula-proposed-clinerules}
Because this review was
restricted to \texttt{docs/}, the following content is the canonical
proposed implementation-agent rule file. It should be materialised at
\texttt{src/training/.clinerules} once source-tree edits are allowed.

\begin{lstlisting}[basicstyle=\ttfamily\scriptsize,breaklines=true]
# =============================================================================
# .clinerules - Simula implementation and maintenance agent
# =============================================================================

Mission:
- Implement and maintain the Simula training data framework in src/training.
- Treat the Simula PDF spec as normative over legacy pipeline behaviour.

Repository reality check:
- The canonical agent location is src/training.
- The default runtime path is still the legacy CSV/template pipeline.
- Simula modules already exist in src/training/pipeline and must become the
  authoritative execution path.

Read these files first on every Simula task:
- docs/pdf/specs/simula-training-spec.pdf
- docs/latex/specs/simula/chapters/12-implementation-instructions.tex
- src/training/scripts/run_hana_pipeline.py
- src/training/pipeline/main.py
- src/training/pipeline/Makefile
- src/training/pipeline/simula_data_generator.py
- src/training/pipeline/simula_taxonomy_builder.py
- src/training/pipeline/hana_schema_extractor.py
- src/training/pipeline/simula_critic_validator.py
- src/training/pipeline/simula_coverage_evaluator.py

Non-negotiable implementation rules:
1. Promote the Simula HANA-direct path over the legacy template path.
2. Keep a single canonical data contract across TrainingExample, Taxonomy,
   evaluation reports, API payloads, and UI artefacts.
3. Do not add new evaluator fields in one module without updating every
   producer and consumer.
4. Wire evaluation into the runtime path:
   - complexity calibration
   - critic calibration
   - taxonomy quality
   - coverage
   - diversity
   - student-teacher gap reporting
5. Emit stable artefacts for every run:
   - extracted_schema.json
   - taxonomies.json
   - training_examples.jsonl (or one agreed canonical JSONL name)
   - pipeline_stats.json
   - evaluation reports
6. Make API and frontend consume the same artefact contract as the runner.
7. Preserve legacy pipeline code only as an explicit compatibility path.

Required tests before considering a change complete:
- unit tests for any changed Simula module
- contract test for TrainingExample and Taxonomy serialization
- integration test for the runner entrypoint used by the API
- regression test proving frontend/API artefact paths still resolve

Definition of done:
- the default pipeline command launches the Simula flow
- evaluator modules execute without contract errors
- governance captures run status and evaluation outputs
- frontend surfaces the real Simula artefacts instead of legacy placeholders
\end{lstlisting}

\section{Proposed runtime-monitor agent rule file}
\label{sec:simula-proposed-monitor-rules}
The runtime-monitor agent below is the canonical proposed companion to
the implementation agent. It should be materialised later at a path
such as \texttt{src/training/.clinerules.runtime-monitor} or another
agreed operational-agent location.

\begin{lstlisting}[basicstyle=\ttfamily\scriptsize,breaklines=true]
# =============================================================================
# .clinerules - Simula runtime monitor and alerting agent
# =============================================================================

Mission:
- Monitor live Simula runs and raise actionable alerts.
- Distinguish live production runs from preview/demo fallbacks.

Primary evidence sources:
- governance training runs and gate checks
- /pipeline/status and websocket pipeline logs
- /jobs and job websocket updates
- HANA health and mode (live vs preview)
- vLLM / TurboQuant health
- emitted run artefacts in the agreed output directory

Treat these as critical alert conditions:
1. Pipeline state = error or job status = failed.
2. HANA extraction completed with zero tables.
3. Taxonomy build completed with zero taxonomies.
4. Example generation completed with zero examples.
5. Expected artefact missing:
   - extracted_schema.json
   - taxonomies.json
   - training JSONL
   - pipeline_stats.json
6. Any evaluator raises a contract error or runtime exception.
7. Production-governed run is operating in preview mode.
8. Coverage or taxonomy quality falls below agreed thresholds.
9. Critic calibration degrades:
   - p(accept | correct) <= 0.80
   - p(accept | corrupted) >= 0.30
10. Diversity falls outside target band:
   - global diversity materially below target
   - local diversity materially below target

Recommended severities:
- Critical:
  failed run, preview mode on a production run, missing core artefact,
  contract drift in evaluator modules
- High:
  zero extracted tables, zero taxonomies, zero examples,
  critic calibration failure
- Medium:
  low coverage, low diversity, repeated governance blocks

Required alert payload:
- run_id
- job_id
- team or tenant
- stage
- severity
- exact failing condition
- last relevant log lines
- affected artefact path
- recommended next action

Response rules:
1. Open a notification and attach it to the owning user/team.
2. Append a governance note for Critical and High alerts.
3. Never report "healthy" if the run is only in preview mode.
4. Prefer precise evidence over generic summaries.
5. Collapse duplicate alerts for the same run and condition.
\end{lstlisting}

\section{Sequenced remediation order}
The findings above should be addressed in the following order:
\begin{enumerate}[label=\arabic*.]
  \item Replace the default execution path and API worker wiring so
    the Simula runner becomes authoritative.
  \item Repair the \texttt{TrainingExample}/\texttt{Taxonomy}/evaluator
    contract drift.
  \item Normalise run artefact naming and paths across runner, API,
    governance, and frontend.
  \item Integrate the evaluation framework into the operational run.
  \item Replace simulated UI telemetry with real runtime metrics and
    evaluator outputs.
  \item Materialise the two proposed agent rule files once source-tree
    edits outside \texttt{docs/} are permitted.
\end{enumerate}
